% Complete independent proof fragment. Requires amsmath, amssymb, amsthm,
% amscd, hyperref and theorem/proposition/lemma environments. Prefix ACF.
\section{The absolute arithmetic curve over the full support spectrum}
\subsection{Sources, conventions, and the objects retained}
Alain Connes and Caterina Consani's original author source
\nolinkurl{FF.tex}, \emph{On the Absolute Geometry of Spec Z and the
Fargues--Fontaine curve},
\href{https://arxiv.org/abs/2606.06604}{arXiv:2606.06604}, was read in
full, lines 1--1764. Its SHA-256 is
\begin{center}\small\texttt{14c17a95040be7bfc57bb5f87e58907a05}\\
\texttt{cacd0bcb1b6967b2892cf3d9bc1229}.\end{center}
The precise inputs are Section \texttt{sect9.1}, Proposition
\texttt{prop12}, and the following source locators:
\begin{center}\small
\begin{tabular}{l}
\texttt{prop:monoid\_structure\_morphisms}\\
\texttt{defnloc},\quad \texttt{propflat}\\
\texttt{prop:archimedean\_points\_C}\\
\texttt{prop:nonarchimedean\_points\_C}.
\end{tabular}
\end{center}
The notation of that source is
preserved:
\begin{equation}\tag{ACF1}
 \theta(\eta)=[\mathbb Z],\qquad
 \Theta(\eta)=[\{0\}],\qquad
 \Theta(p)=[\mathbb Z[1/p]],\qquad \kappa\Theta=\theta.
\end{equation}
The first equality concerns an arithmetic divisor class; the second
concerns a point of the topos. In particular the source sheaf is
\begin{equation}\tag{ACF2}
 \begin{gathered}
 \mathcal F=\Theta^{-1}\mathcal O_{\mathbf F_1},\qquad
 \mathcal O_{\mathbf F_1}=\mathbf F_1[T],\\
 \mathcal F_\eta=\mathbf F_1,\quad
 \mathcal F_p=\mathbf F_1[T^{H_p^+}],\quad
 H_p=\mathbb Z[1/p],\quad H_p^+=H_p\cap[0,\infty).
 \end{gathered}
\end{equation}
The original sheaf on the topos and its generic stalk are different
objects in this formula.

For the coefficient foundation the input is \emph{Split Support
Geometry: Universal Zero Fibres, Arithmetic Curves, and Frobenius-Perfect
Quantization}, source byline \emph{The Clankers}, June 2026,
pre-publication Version 11. The complete original proof file is
\href{https://github.com/KokunoYumeto/zeta-function-research-reader/blob/a2851f9eb352e7e4376bc629de86fa4ed9c1c434/workbenches/splitzero-tandem/foundations/split-support-geometry-v11/split_support_geometry_arithmetic_curve_v11.tex}{pinned at commit a2851f9e}.
The locators used here are Definition \texttt{def:\allowbreak lattice-\allowbreak split},
Theorems \texttt{thm:\allowbreak lattice-\allowbreak ideal-\allowbreak classification},
\texttt{thm:\allowbreak lattice-\allowbreak prime-\allowbreak classification},
\texttt{thm:\allowbreak spectral-\allowbreak ordinal-\allowbreak sum}, Proposition
\texttt{prop:\allowbreak classical-\allowbreak collapse}, and Definition
\texttt{def:\allowbreak split-\allowbreak absolute-\allowbreak point-\allowbreak map}. Their relevant source ranges
1221--1311, 1606--1749, and 4591--4788 were read. All extensions proved
below retain every lattice label and reproduce the necessary arguments.

Let $L$ be a bounded distributive lattice, including the possible
one-element lattice. For a nonzero commutative ring $R$ set
\begin{equation}\tag{ACF3}
 \begin{gathered}
 S_R=G_L(R)=\{z_\lambda=(0,\lambda):\lambda\in L\}
       \cup\{\widehat a=(a,1_L):a\in R\},\\
 (a,\lambda)+(b,\mu)=(a+b,\lambda\vee\mu),\qquad
 (a,\lambda)(b,\mu)=(ab,\lambda\wedge\mu),\\
 \tau=z_{0_L},\quad e=z_{1_L}=\widehat0,\quad
 1=\widehat1,\qquad q(a,\lambda)=a.
 \end{gathered}
\end{equation}
These operations close on the displayed carrier: an amplitude different
from zero forces support $1_L$, while a factor with smaller support has
zero amplitude. The ring and lattice distributive laws give the semiring
laws component by component. Thus $q$ is a unital semiring map. The
notation $\widehat0$ always denotes $e$, not $\tau$.

\subsection{The full prime spectrum and its arithmetic retraction}
Here $\operatorname{Spec}$ means prime ideals of the additive semiring:
an ideal contains $\tau$, is closed under addition, and absorbs
multiplication by every element. It does not mean prime ideals of only
the underlying multiplicative monoid. Put $S=S_{\mathbb Z}$,
$X_L=\operatorname{Spec}S$, and $Y=\operatorname{Spec}\mathbb Z$.

\begin{proposition}[All primes and all basic opens]
Write $\operatorname{Spec}_{\rm lat}L$ for the proper prime lattice ideals,
where ideals are lower, contain $0_L$, and are closed under finite joins.
The points of $X_L$ are exactly
\begin{equation}\tag{ACF4}
 \begin{aligned}
 P_J&=\{z_\lambda:\lambda\in J\},
                  &&J\in\operatorname{Spec}_{\rm lat}L,\\
 Q_{\mathfrak p}&=q^{-1}(\mathfrak p),
                  &&\mathfrak p\in Y.
 \end{aligned}
\end{equation}
Every $P_J$ is strictly contained in $Q_{(0)}$, and
$Q_{(0)}\subsetneq Q_{(p)}$ for every rational prime $p$.
Within the support stratum inclusion of primes is exactly inclusion of
$J$. With $A=\operatorname{Spec}_{\rm lat}L$, the basic opens are
\begin{equation}\tag{ACF5}
 \begin{aligned}
 D_S(z_\lambda)&=\{P_J:\lambda\notin J\},\\
 D_S(\widehat n)&=A\ \cup\ i(D_{\mathbb Z}(n))\quad(n\in\mathbb Z),\\
 D_S(e)&=A,\qquad V_S(e)=i(Y),\qquad
 i(\mathfrak p)=Q_{\mathfrak p}.
 \end{aligned}
\end{equation}
In the second line $n=0$ is allowed, and then the right side is $A$.
\end{proposition}
\begin{proof}
If an ideal $I$ contains $\widehat n$ with $n\ne0$, then multiplication
by $\widehat{-1}$ and addition give $e\in I$. If it contains $e$
initially the same conclusion is already available. Then
$ez_\lambda=z_\lambda$ puts the entire zero fibre in $I$.
The amplitudes whose supported representatives belong to $I$ form a ring
ideal, so $I=q^{-1}(\mathfrak a)$ for a unique ring ideal
$\mathfrak a$. Otherwise $I$ consists solely of $z_\lambda$ for a proper
lattice ideal $J$: addition gives join closure and
$z_\mu z_\lambda=z_\mu$ when $\mu\le\lambda$ gives downward closure.
Conversely these constructions are ideals. The amplitude ideals are
prime exactly when $\mathfrak a$ is prime, by surjectivity of $q$.
For a support ideal, a product of two top-supported elements is
still top-supported, even if its amplitude is zero. A product
of a top-supported element with $z_\lambda$ has support $\lambda$.
For two zero elements primality is precisely
$\lambda\wedge\mu\in J\Rightarrow\lambda\in J$ or $\mu\in J$.
This proves exhaustion and primality. The strict cross-stratum inclusion
follows because every $Q_{\mathfrak p}$ contains $e$ and no $P_J$ does.
Membership in the displayed ideals proves every formula for the basic
opens directly.
\end{proof}

\begin{proposition}[A retraction which retains the support stratum]
The map $i:Y\to X_L$ is a homeomorphism onto the closed subspace $V(e)$.
There is a continuous retraction
\begin{equation}\tag{ACF6}
 \begin{gathered}
 r:X_L\longrightarrow Y,\qquad
 r(P_J)=\eta,\quad r(Q_{(0)})=\eta,\quad r(Q_{(p)})=(p),\\
 ri=\operatorname{id}_Y.
 \end{gathered}
\end{equation}
Its fibres are $r^{-1}(\eta)=A\cup\{Q_{(0)}\}$ and
$r^{-1}((p))=\{Q_{(p)}\}$. For $n\ne0$,
\begin{equation}\tag{ACF7}
 r^{-1}D_{\mathbb Z}(n)=D_S(\widehat n).
\end{equation}
For $n=0$ the left side is empty, whereas $D_S(\widehat0)=A$.
\end{proposition}
\begin{proof}
The formulas in ACF5 give the homeomorphism onto $V(e)$ and show that
ACF7 holds for every nonzero integer, since every nonempty basic open of
$Y$ contains $\eta$. The inverse image of the empty open is empty.
These opens form a basis, proving continuity. All statements about fibres
and the composite follow from the explicit assignments. In particular
$r$ is a quotient map: if $r^{-1}U$ is open, then
$U=i^{-1}r^{-1}U$ is open.
\end{proof}
For $0_L\ne1_L$, this continuous retraction is not induced by a unital
semiring homomorphism $\mathbb Z\to S$. Indeed any such homomorphism
would send $1$ to $\widehat1$. If $u$ were the image of $-1$, then
$\widehat1+u=\tau$ would be required. Its left support is $1_L$, which
cannot equal $0_L$. The continuous map and its sheaf pullback nevertheless
exist exactly as proved.

\subsection{The pullback sheaf and its exact comparison}
Define a spherical-algebra sheaf on the actual full base by
\begin{equation}\tag{ACF8}
 \mathcal F_L=r^{-1}\mathcal F=(\Theta\circ r)^{-1}
                           \mathcal O_{\mathbf F_1}.
\end{equation}
For a continuous map, the inverse image is the sheafification of the
presheaf assigning to $W\subseteq X_L$ the colimit of $\mathcal F(V)$
over opens $V\subseteq Y$ for which $W\subseteq r^{-1}V$.
This specifies the sheaf, not only a list of stalks.

\begin{proposition}[Exact restriction, stalks, and direct image]
There are canonical isomorphisms
\begin{equation}\tag{ACF9}
 i^{-1}\mathcal F_L\simeq\mathcal F,\qquad
 r_*\mathcal F_L\simeq\mathcal F,\qquad
 (\mathcal F_L)_x\simeq\mathcal F_{r(x)}.
\end{equation}
Consequently
\begin{equation}\tag{ACF10}
 (\mathcal F_L)_{P_J}=\mathbf F_1,\qquad
 (\mathcal F_L)_{Q_{(0)}}=\mathbf F_1,\qquad
 (\mathcal F_L)_{Q_{(p)}}=\mathbf F_1[T^{H_p^+}].
\end{equation}
All generic-direction stalk maps from the last object to either of the
first two are
\begin{equation}\tag{ACF11}
 \epsilon_p(0)=0,\qquad \epsilon_p(T^h)=1\quad(h\in H_p^+).
\end{equation}
Maps between stalks $\mathbf F_1$ in ACF10 are identities.
\end{proposition}
\begin{proof}
The stalk of the inverse-image presheaf at $x$ is the colimit of
$\mathcal F(V)$ over neighbourhoods of $r(x)$: every such $V$ yields the
neighbourhood $r^{-1}V$ of $x$, and every representative of the presheaf
stalk is already represented on some such $V$. The two identifications
respect equality of germs. Sheafification preserves this stalk, proving
the third isomorphism. The first follows either by the same construction
and $ri=\operatorname{id}$ or by the canonical composition isomorphism
for inverse images.

The second isomorphism holds for any sheaf $E$ on $Y$ in place of
$\mathcal F$. The adjunction unit sends $E(U)$ into
$(r^{-1}E)(r^{-1}U)$. Restriction along $i$ gives a left inverse using
$i^{-1}r^{-1}E=E$. If two sections over $r^{-1}U$ have the same restriction,
their germs agree at every point $Q_{\mathfrak p}$ with $\mathfrak p\in U$.
When $U$ is nonempty it contains $\eta$. Each support point is a
generization of $Q_{(0)}$, so the germ at that support point is the
image of the germ at $Q_{(0)}$. The two sections therefore agree at all
stalks and are equal. The case $U=\varnothing$ is automatic. This proves
bijectivity, respecting all sheaf operations and restrictions.

The source has all endomorphisms
$\operatorname{Fr}_m(T^h)=T^{mh}$, including $m=0$, with the absorbing
zero fixed. Every sheaf generization map commutes with them. On
$\mathbf F_1$, $\operatorname{Fr}_0$ is the identity. Thus a generization
$g$ to $\mathbf F_1$ satisfies
$g(T^h)=\operatorname{Fr}_0g(T^h)=g(T^0)=1$.
It sends the absorbing zero to zero. This proves ACF11 without selecting
an unsupported alternative augmentation. All other displayed maps are
unital maps $\mathbf F_1\to\mathbf F_1$, hence identities.
\end{proof}

\begin{proposition}[The actual semiring structure sheaf]
Let $\mathcal O_X$ be the sheaf of locally represented fractions of $S$
on $X_L$. For a lattice prime $J$ define
\begin{equation}\tag{ACF12}
 \begin{gathered}
 L_J=L[(L\setminus J)^{-1}],\\
 [\lambda]=[\mu]\ \Longleftrightarrow\
 \text{there is }\nu\notin J\text{ with }
                  \lambda\wedge\nu=\mu\wedge\nu.
 \end{gathered}
\end{equation}
Its operations are induced join and meet. The stalks are
\begin{equation}\tag{ACF13}
 (\mathcal O_X)_{P_J}=L_J,\qquad
 (\mathcal O_X)_{Q_{(0)}}=G_L(\mathbb Q),\qquad
 (\mathcal O_X)_{Q_{(p)}}=G_L(\mathbb Z_{(p)}).
\end{equation}
The generic-direction map from an arithmetic stalk to $L_J$ is
$(a,\lambda)\mapsto[\lambda]$. The maps between arithmetic stalks
are the usual localization of amplitudes with support unchanged.
If $J\subseteq K$, the map $L_K\to L_J$ is the canonical localization.
\end{proposition}
\begin{proof}
The locally represented fraction definition gives the stalk $S_P$ at
a prime $P$: a fraction near $P$ has denominator outside $P$, and equality
of fraction germs is the localization equality witnessed by another
such denominator. At $Q_{\mathfrak p}$, every allowed denominator is
$\widehat n$ with $n\notin\mathfrak p$ and $n\ne0$. Thus
\[
 (a,\lambda)/\widehat n\longmapsto(a/n,\lambda)
\]
is onto $G_L(\mathbb Z_{\mathfrak p})$. The usual localization
cross-multiplication criterion proves injectivity: an additional
denominator is top-supported, so equality preserves the supports
exactly, and its amplitude is a nonzero integer. Operations match
component by component.

At $P_J$, $e$ is an allowed denominator. The relation $e^2=e$ makes $e=1$
after localization. Hence every $(a,\lambda)$ equals its product with
$e$, namely $z_\lambda$. This first localization is exactly $L$:
$z_\lambda\mapsto\lambda$ is the isomorphism. The remaining allowed
zero denominators are the $z_\nu$ with $\nu\notin J$, so the second
localization is $L_J$. An invertible idempotent is $1$; thus every
fraction is represented by its numerator, and the usual fraction
equality becomes precisely ACF12. This also proves the formula for
all maps, because they are induced by the original localization maps.
The top and bottom of $L_J$ remain different: equality would require
$\nu\notin J$ with $\nu=0_L$, impossible for a lattice ideal.
\end{proof}

\begin{proposition}[Coefficient comparison over the whole base]
On $Y$ define the sheaf $B_L$ by
$B_L(U)=G_L(\mathcal O_Y(U))$ for nonempty opens $U$, and the terminal
semiring for the empty open. There is a canonical sheaf morphism
\begin{equation}\tag{ACF14}
 \Phi:r^{-1}B_L\longrightarrow\mathcal O_X
\end{equation}
which is an isomorphism at every arithmetic stalk and a surjection
\begin{equation}\tag{ACF15}
 G_L(\mathbb Q)\longrightarrow L_J,\qquad
 (a,\lambda)\longmapsto[\lambda]
\end{equation}
at $P_J$. Its exact equality relation at that stalk is equality of the
classes $[\lambda]$ and $[\mu]$ in ACF12, regardless of amplitudes.
\end{proposition}
\begin{proof}
Every two nonempty opens of $Y$ meet, and their rings of sections embed
in $\mathbb Q$. Compatible local sections therefore have one support
label on all members of a nonempty cover and glue their amplitudes by
the ring sheaf axiom. A nonzero glued amplitude forces top support.
This proves the claimed sheaf description. Its stalks are
$G_L(\mathbb Z_{(p)})$ and $G_L(\mathbb Q)$: passage to the stalk keeps
every support label and takes the corresponding filtered localization
of the amplitude.

For $D_{\mathbb Z}(n)$ with $n\ne0$, a section of $B_L$ is represented
by $(a/n^k,\lambda)$. Send it to the local fraction
$(a,\lambda)/\widehat{n^k}$ over $D_S(\widehat n)$, which is its full
inverse image by ACF7. Equality and restriction are preserved by the
fraction rules. These maps on a basis give $B_L\to r_*\mathcal O_X$
and hence ACF14. At arithmetic stalks the map is the isomorphism just
proved. At a support stalk all nonzero integers have been inverted,
and the subsequent inversion of $e$ and $z_\nu$ gives ACF15. Every
$[\lambda]$ is hit by $z_\lambda$. Formula ACF12 gives precisely its
fibres, proving all assertions.
\end{proof}
The functor used here is Connes--Consani's $H$ from
\href{https://arxiv.org/abs/1502.05585}{\emph{Absolute algebra and
Segal's $\Gamma$-rings: au dessous de $\overline{\operatorname{Spec}
\mathbb Z}$}}, original \nolinkurl{F1_corr.tex}, formula
\texttt{functsum} (lines 433--442) and Lemma \texttt{sssalg}
(lines 519--533). That source was read completely, lines 1--1336,
for the preceding coefficient derivation.
Applying the semiring functor $H$ degree by degree gives the corresponding
spherical-algebra map $r^{-1}HB_L\to H\mathcal O_X$. Here
$(HA)(k_+)=A^k$, a pointed map sums coordinates over its fibres, and
multiplication sends $(a_i),(b_j)$ to $(a_i b_j)_{i,j}$. These formulas
show directly that the comparison commutes with all pointed maps and
multiplications. Stalks and inverse images commute with this finite
coordinate construction. Thus ACF14 is also the exact connection with
the split coefficient sheaf, rather than an identification with
$\mathcal F_L$.

\begin{proposition}[The spherical comparison is forced by the generic stalk]
There is exactly one unital spherical-algebra sheaf map
\begin{equation}\tag{ACF16}
 \mathcal F_L\longrightarrow H\mathcal O_X.
\end{equation}
At $Q_{(p)}$ it sends every $T^h$ to $1$ and the source absorbing zero
to the target additive zero. At all other stalks it is the unit map from
$\mathbf F_1$.
\end{proposition}
\begin{proof}
The original endomorphism $\operatorname{Fr}_0$ of
$\mathcal O_{\mathbf F_1}=\mathbf F_1[T]$ factors through the constant
spherical algebra $\mathbf F_1$: every monomial goes to $1$ and the
absorbing zero to $0$. This factorization is equivariant for every
$\operatorname{Fr}_m$, so it pulls back to a sheaf augmentation
$\mathcal F_L\to\mathbf F_1$. Compose with the unit map to
$H\mathcal O_X$ to obtain ACF16.
For uniqueness, any such morphism has the unique unit map at $Q_{(0)}$.
The generization $G_L(\mathbb Z_{(p)})\to G_L(\mathbb Q)$ is injective.
Compatibility with ACF11 forces the image of every $T^h$ at $Q_{(p)}$
to be the unique element whose image in $G_L(\mathbb Q)$ is $1$,
namely $1$. At support points the source is $\mathbf F_1$, whose unital
map is unique. A spherical monoid-algebra map is determined by its
values on monomials, as is also proved explicitly below. Hence the
maps agree on every stalk and are equal as sheaf maps.
\end{proof}

\subsection{Every point of the finite-prime stalk, with every zero label}
Put $S_{\mathbb C}=G_L(\mathbb C)$. A spherical-algebra point of
$\mathbf F_1[T^{H_p^+}]$ with values in $HS_{\mathbb C}$ is exactly a
pointed unital multiplicative monoid map to $S_{\mathbb C}$.
Explicitly, for a pointed monoid $M$, the source in degree $k_+$ is
$M\wedge k_+$. A monoid map $f:M\to S_{\mathbb C}$ induces the map
which sends $(m,j)$ to the tuple with $f(m)$ in coordinate $j$ and
$\tau$ elsewhere. The wedge basepoint goes to the all-$\tau$ tuple.
A pointed map between finite pointed sets moves this one occupied
coordinate or sends it to the basepoint; the coordinate-sum formula of
$H$ does exactly the same. Products agree because $f$ is multiplicative.
Conversely the maps $1_+\to k_+$ selecting one coordinate force this
formula for every natural map. Its degree-one component must preserve
multiplication, the unit and the absorbing zero. This proves the
source adjunction in the present case, including its zero convention.

Consequently a point is a character
\begin{equation}\tag{ACF17}
 \chi:H_p^+\longrightarrow S_{\mathbb C},\qquad
 \chi(0)=1,\quad \chi(h+h')=\chi(h)\chi(h'),
\end{equation}
with the additional source absorbing zero mapped to $\tau$.
The exponent $0$ and the absorbing zero are distinct source elements.

\begin{proposition}[Complete root-sequence classification]
Define
\begin{equation}\tag{ACF18}
 \mathcal U_p=
 \{(u_n)_{n\ge0}\in(\mathbb C^\times)^{\mathbb N}:
                                      u_{n+1}^p=u_n\}.
\end{equation}
Then the full monoid of points is the disjoint union
\begin{equation}\tag{ACF19}
 \begin{gathered}
 \mathcal P_{p,L}=\mathcal U_p\ \sqcup\ \{Z_\lambda:\lambda\in L\},\\
 Z_\lambda(0)=1,\qquad Z_\lambda(h)=z_\lambda\quad(h>0).
 \end{gathered}
\end{equation}
A sequence $u\in\mathcal U_p$ gives the character
\begin{equation}\tag{ACF20}
 \chi_u(a/p^n)=\widehat{u_n^a}\quad(a\in\mathbb N),
 \qquad \chi_u(0)=1.
\end{equation}
Products and the amplitude map are exactly
\begin{equation}\tag{ACF21}
 \begin{gathered}
 \chi_u\chi_v=\chi_{uv},\qquad
 \chi_u Z_\lambda=Z_\lambda,\qquad
 Z_\lambda Z_\mu=Z_{\lambda\wedge\mu},\\
 q_*:\mathcal P_{p,L}\longrightarrow
       \varprojlim_{x\mapsto x^p}\mathbb C
       =\mathcal U_p\sqcup\{\boldsymbol0\},\\
 q_*^{-1}(\boldsymbol0)=\{Z_\lambda:\lambda\in L\},\qquad
 q_*^{-1}(u)=\{\chi_u\}\quad(u\in\mathcal U_p).
 \end{gathered}
\end{equation}
This inverse limit is asserted as a multiplicative monoid, not as a
perfectoid field or as the ordinary field $\mathbb C$.
\end{proposition}
\begin{proof}
Given a character put $s_n=\chi(1/p^n)$. Then $s_{n+1}^p=s_n$.
Conversely, from any sequence with this property define
$\chi(a/p^n)=s_n^a$, with $s_n^0=1$. To check independence of the
representation, pass to a common denominator $p^N$:
$s_n^a=s_N^{a p^{N-n}}$; equal rational exponents give equal integer
powers. The same common-denominator calculation proves multiplicativity.

Write $s_n=(u_n,\lambda_n)$. Idempotence of lattice meet gives
$\lambda_{n+1}=\lambda_n$ for every $n$. If one $u_n$ is zero, the
relations and the absence of nonzero nilpotents in $\mathbb C$ force
all $u_n=0$. Thus all $s_n=z_\lambda$ for a single arbitrary label
$\lambda$, producing $Z_\lambda$. Otherwise all $u_n$ are nonzero
and every support is $1_L$, producing exactly ACF18 and ACF20.
All products follow by evaluating at the generators $1/p^n$.
The unit is the constant character $1$ in $\mathcal U_p$, and the
absorbing point is $Z_{0_L}$. The amplitude formula follows coordinate
by coordinate; a complex coherent sequence is either everywhere zero
or everywhere nonzero. This proves exhaustion, injectivity, and every
fibre statement.
\end{proof}
The labelled zeros form a multiplicatively absorbing ideal of the
point monoid, with multiplication given by meet. They remain distinct
before the explicitly stated amplitude quotient. Each support prime
$P_J$ has a unique algebraic $HS_{\mathbb C}$-valued stalk point, because
its stalk is $\mathbf F_1$; these points are still indexed by their
different base points $P_J$. The map on point sets induced by the
stalk augmentation ACF11 sends this unique generic point to the
constant unit character in ACF19. It does not select a zero branch.

\begin{proposition}[Every Frobenius, including exponent zero]
For every integer $m\ge1$, source Frobenius induces
\begin{equation}\tag{ACF22}
 \operatorname{Fr}_m^*(\chi_u)=\chi_{(u_n^m)_n},\qquad
 \operatorname{Fr}_m^*(Z_\lambda)=Z_\lambda.
\end{equation}
For $m=0$, every point is sent to the constant unit character.
The $p$-th Frobenius is a bijection, with inverse on the group part
$(u_n)\mapsto(u_{n+1})$, and with every $Z_\lambda$ fixed.
More generally, write $m=p^a d$ with $(d,p)=1$. On $\mathcal U_p$,
$\operatorname{Fr}_m^*$ is surjective with kernel canonically isomorphic
to $\mu_d(\mathbb C)$ by the zeroth coordinate after removal of the
$p$-power automorphism. For $d>1$ it is not injective.
All these maps commute with the sheaf generizations ACF11.
\end{proposition}
\begin{proof}
Precomposition sends $\chi(h)$ to $\chi(mh)=\chi(h)^m$, giving ACF22.
For $m=0$ it gives $\chi(0)=1$ at every exponent, while the extra source
absorbing zero is still mapped to $\tau$. The inverse for $p$ follows
from $u_{n+1}^p=u_n$ and fixes the constant support sequences.
For $d$ prime to $p$, to lift a sequence $u$ choose $v_0^d=u_0$.
Having chosen $v_n$, choose a $p$-th root $w$ of $v_n$.
Then $w^d/u_{n+1}$ is a $p$-th root of unity. Since raising to the
$d$-th power permutes $\mu_p$, multiply $w$ by the unique suitable
$p$-th root of unity to obtain $v_{n+1}$ with both
$v_{n+1}^p=v_n$ and $v_{n+1}^d=u_{n+1}$. This constructs a preimage.
In the kernel every coordinate belongs to $\mu_d$, and raising to
$p$ is an automorphism of that finite group, so $v_0$ determines the
whole sequence uniquely. Combining with the $p$-power bijection proves
the assertion for $m$. Finally $\epsilon_p(T^{mh})=1$ for every $h,m$,
which proves the compatibility even at $m=0$.
\end{proof}

\subsection{The source's actual locality condition}
In \texttt{defnloc}, Connes--Consani define a local character using a
homomorphism from the ordered group $H$ into a complete Hausdorff
topological \emph{group}. It must be continuous for the topology induced
by its place and then extends uniquely to the corresponding local field.
For $H_p$ the places are $p$ and $\infty$; for the trivial group there
are no such places (source Proposition \texttt{classpla}). We apply
this definition exactly to the unit group of the present target.

\begin{proposition}[The exact boundary excluded by group extension]
The units of $S_{\mathbb C}$ are exactly
$\{\widehat a:a\in\mathbb C^\times\}$, identified with the usual
topological group $\mathbb C^\times$. A character ACF17 extends to a
group homomorphism $H_p\to S_{\mathbb C}^{\times}$ if and only if it
belongs to $\mathcal U_p$. Every $Z_\lambda$ fails this extension,
including $Z_{1_L}$ and $Z_{0_L}$.
\end{proposition}
\begin{proof}
If $(a,\lambda)(b,\mu)=1$, then $ab=1$, so $a,b$ are nonzero and both
supports are $1_L$. Conversely $\widehat a^{-1}=\widehat{a^{-1}}$.
A character extending to a group must take every $h>0$ to a unit,
since its value at $-h$ is an inverse. This excludes every zero branch.
A character from $\mathcal U_p$ extends by
$\rho(h-k)=\chi(h)\chi(k)^{-1}$ for $h,k\in H_p^+$;
if $h-k=h'-k'$, then $h+k'=h'+k$, proving independence of the
representation. The group law and uniqueness follow from the same
calculation. Thus the obstruction object is exactly the retained
lattice ideal $\{Z_\lambda\}$ of ACF21.
\end{proof}
This argument does not discard the algebraic boundary: it identifies
its exact failure to satisfy the source's group-valued definition.
For an explicit topological semiring realization give $L$ the
discrete topology and $S_{\mathbb C}\subseteq\mathbb C\times L$ the
subspace topology. Its operations are continuous, since the ring
operations are continuous and maps between finite products of
discrete sets are continuous. Its unit group has exactly the stated
topology. Even without requiring group extension, a zero-branch
character is not continuous at exponent $0$: at the archimedean place
$1/p^n\to0$, and at the $p$-adic place $p^n\to0$, whereas its amplitudes
are constantly $0$ along these sequences and its amplitude at $0$ is
$1$.

\begin{proposition}[All archimedean local characters]
The archimedean local points form the family
\begin{equation}\tag{ACF23}
 \chi_z(h)=\widehat{\exp(zh)},\qquad z\in\mathbb C,\quad h\in H_p^+.
\end{equation}
Different $z$ give different characters. The trivial point is $z=0$.
Source Frobenius $\operatorname{Fr}_m^*$ acts by $z\mapsto mz$.
\end{proposition}
\begin{proof}
The group $H_p$ is dense in $\mathbb R$: rounding $p^n t$ to the nearest
integer approximates any real $t$ with error at most $1/(2p^n)$.
A continuous homomorphism $H_p\to\mathbb C^\times$ is uniformly
continuous for the group uniformities. Completeness of
$\mathbb C^\times$ for its group uniformity follows from
$\mathbb C^\times\simeq\mathbb R\times S^1$, by logarithm of the modulus
and phase. Therefore the homomorphism extends uniquely to $\mathbb R$:
apply it to approximating Cauchy sequences, and use uniform continuity
to prove independence and continuity of the limit.
For a continuous homomorphism $f:\mathbb R\to\mathbb C^\times$, choose a
small interval around zero on which $f$ takes values in a neighbourhood
of $1$ with a continuous logarithm $g$ and $g(0)=0$.
For sufficiently small $s,t,s+t$, the difference
$g(s+t)-g(s)-g(t)$ belongs to $2\pi i\mathbb Z$ and is continuous;
on the connected neighbourhood of $(0,0)$ it is therefore zero.
Continuous local additivity gives $g(t)=zt$: rational subdivisions
first give this equality on rational multiples of a fixed sufficiently
small nonzero real argument, and continuity gives it on the entire
small interval. Subdividing any real $t$ into sufficiently small equal
pieces gives $f(t)=\exp(zt)$ everywhere. Conversely each such function
is continuous and multiplicative. If two parameters give the same
character on $H_p$, density makes them agree on $\mathbb R$;
their local logarithms have equal slopes. Frobenius follows by
substituting $mh$.
\end{proof}

For the $p$-adic family fix its parameter convention explicitly.
For $x\in\mathbb Q_p$, let $\{x\}_p\in\mathbb Z[1/p]\cap[0,1)$ be its
finite sum of negative $p$-adic digits. Its class modulo $\mathbb Z$
depends only on $x\bmod\mathbb Z_p$. Set
\begin{equation}\tag{ACF24}
 \psi_p(x)=\exp(2\pi i\{x\}_p).
\end{equation}
The map $\mathbb Z[1/p]/\mathbb Z\to\mathbb Q_p/\mathbb Z_p$ is a
bijection: every $p$-adic class has the finite negative-digit
representative and its kernel is
$\mathbb Z[1/p]\cap\mathbb Z_p=\mathbb Z$.
It respects addition, so ACF24 is a continuous additive character
with kernel $\mathbb Z_p$.

\begin{proposition}[All $p$-adic local characters]
The $p$-adic local points form the family
\begin{equation}\tag{ACF25}
 \chi_a(h)=\widehat{\psi_p(ah)},\qquad a\in\mathbb Q_p,
 \quad h\in H_p^+.
\end{equation}
The parameter is unique, the trivial point is $a=0$, and source
Frobenius acts by $a\mapsto ma$.
\end{proposition}
\begin{proof}
The group $H_p$ is dense in $\mathbb Q_p$: truncate the $p$-adic
expansion of any number at successively larger nonnegative exponents.
The completeness argument just used extends each continuous character
uniquely to $\mathbb Q_p$. A continuous homomorphism
$f:\mathbb Q_p\to\mathbb C^\times$ maps a sufficiently small open
subgroup $p^N\mathbb Z_p$ into a neighbourhood of $1$ containing no
nontrivial subgroup. Such a neighbourhood exists: use a bounded open
interval for $\log|z|$ and an open phase arc of length less than one
half-turn; powers of any nonidentity element leave it. Continuity
therefore gives $f(p^N\mathbb Z_p)=1$. For each $x\in\mathbb Q_p$,
some $p$-power multiple belongs to $p^N\mathbb Z_p$, so $f(x)$ is a
root of unity of $p$-power order.

Write $x_n=f(1/p^n)$ and use the fixed isomorphism
$\mathbb Q_p/\mathbb Z_p\to\mu_{p^\infty}(\mathbb C)$ induced by
$\psi_p$ to write $x_n=\psi_p(b_n)$, where $b_n$ is specified modulo
$\mathbb Z_p$. The relation $x_{n+1}^p=x_n$ gives
$p b_{n+1}\equiv b_n\pmod{\mathbb Z_p}$.
Thus $a_n=p^n b_n$ is specified modulo $p^n\mathbb Z_p$ and satisfies
$a_{n+1}\equiv a_n\pmod{p^n\mathbb Z_p}$.
Choose representatives recursively; their differences have valuations
at least $n$, so they form a Cauchy sequence in $\mathbb Q_p$. Its
limit $a$ satisfies $a\equiv a_n\pmod{p^n\mathbb Z_p}$ for every $n$.
It follows that $x_n=\psi_p(a/p^n)$. Agreement on these generators gives
agreement on $H_p$, and continuity gives agreement on $\mathbb Q_p$.
If two parameters work, their difference belongs to
$\bigcap_{n\ge0}p^n\mathbb Z_p=\{0\}$, proving uniqueness.
Conversely ACF24 shows that every $a$ gives a continuous character.
Replacing $h$ by $mh$ proves the Frobenius action.
\end{proof}

\begin{proposition}[The exact intersection of the two local families]
As subsets of the algebraic point set ACF19, their intersection is
\begin{equation}\tag{ACF26}
 \{\chi_{2\pi i b}:b\in\mathbb Z[1/p]\}
  =\{\chi_b\text{ of ACF25}:b\in\mathbb Z[1/p]\}.
\end{equation}
The place-labelled local space therefore has two copies of these
characters when the labels $p$ and $\infty$ are retained. Forgetting
the place label identifies exactly these copies. The zero-labelled
ideal is disjoint from both families.
\end{proposition}
\begin{proof}
If an archimedean character $\exp(zh)$ is also $p$-adic local, then its
value at $h=1$ is a root of unity of $p$-power order by the preceding
proof. Therefore $\exp z=\exp(2\pi i c/p^n)$ for integers $c,n$, so
$z=2\pi i b$ with $b\in\mathbb Z[1/p]$.
For such $b$ and $h\in H_p$, the rational number $bh$ belongs to
$\mathbb Z[1/p]$ and differs from $\{bh\}_p$ by an integer. Hence
$\exp(2\pi i b h)=\psi_p(bh)$. This proves both sufficiency and the
exact match of parameters; uniqueness follows from the preceding
propositions. Exclusion of the zero-labelled ideal was proved by the
unit criterion.
\end{proof}

\subsection{Comparison with the foundation's resolved point diagram}
The foundation's Definition \texttt{def:\allowbreak split-\allowbreak absolute-\allowbreak point-\allowbreak map}
assigns, in the two-element support case, the ordered groups
$\{0\},\mathbb Z,H_p$ to $P_\tau,Q_{(0)},Q_{(p)}$, respectively.
Its Corollary \texttt{cor:\allowbreak resolved-\allowbreak stalk-\allowbreak chain} gives the forward maps
$\{0\}\hookrightarrow\mathbb Z\hookrightarrow H_p$, and hence
$\mathbf F_1\to\mathbf F_1[T]\to\mathbf F_1[T^{H_p^+}]$.
This is a point-and-stalk diagram explicitly described as such in the
source. The sheaf pulled back along $r$ instead has the stalks and the
reverse, generic-direction maps in ACF10--ACF11. In particular it
restricts exactly to CC's original sheaf on $V(e)$; its value at
$Q_{(0)}$ is not $\mathbf F_1[T]$.

There are exact objectwise maps relating the two intermediate objects:
\begin{equation}\tag{ACF27}
 \mathbf F_1\xrightarrow{\text{unit}}\mathbf F_1[T]
 \xrightarrow{T\mapsto1}\mathbf F_1,
 \qquad\text{whose composite is the identity.}
\end{equation}
The following calculation records the precise defect when this
augmentation is compared with the forward inclusion into the finite
stalk. It does not assert a sheaf identification.

\begin{proposition}[The coequalizer of the intermediate comparison]
Let $j,c:\mathbf F_1[T]\rightrightarrows
\mathbf F_1[T^{H_p^+}]$ be the maps $j(T)=T^1$ and $c(T)=1$.
At the level of pointed multiplicative monoids, their coequalizer is
\begin{equation}\tag{ACF28}
 \mathbf F_1[T^{H_p^+}]\longrightarrow
 \mathbf F_1[H_p/\mathbb Z],\qquad
 T^h\longmapsto[h],\quad 0\longmapsto0.
\end{equation}
The target is the group $H_p/\mathbb Z$ with an extra absorbing zero;
its identity is $[0]$, distinct from that zero. Applying the spherical
monoid-algebra functor gives the spherical-algebra coequalizer.
Its $HS_{\mathbb C}$-valued points are exactly the characters ACF25
with $a\in\mathbb Z_p$. They have no extra zero-support branches.
\end{proposition}
\begin{proof}
The imposed relation is $T^1=1$. Every $h=a/p^n\in H_p^+$ then
satisfies $(T^h)^{p^n}=T^a=1$, so all monomials become invertible.
The resulting group is the quotient of the group completion $H_p$ by
$\mathbb Z$. More explicitly, if $h-k\in\mathbb Z$, multiply the
relation $T^1=1$ by the appropriate nonnegative monomial to identify
$T^h$ and $T^k$. Conversely their classes in $H_p/\mathbb Z$ can be
equal only under that condition. The extra absorbing zero remains
separate, witnessed by the displayed map to the pointed group.
Any monoid map equalizing $j,c$ sends $T^1$ to $1$, makes every image
monomial a unit by the displayed power identity, and therefore factors
uniquely through this quotient. This proves the universal property.
The monoid-algebra adjunction transports that same universal property
to spherical algebras: the coordinate-inclusion argument proving it
above applies to every spherical target using its degree-one
multiplication and its structure maps.

A point of this quotient is a coherent sequence of roots of unity
$x_{n+1}^p=x_n$ with $x_0=1$. Since $x_n^{p^n}=1$, the parameter
construction in the proof of ACF25 yields a unique $a\in\mathbb Z_p$;
conversely such $a$ satisfies $\psi_p(a)=1$ and gives a point.
All monomials in a group must map to units, excluding every
$Z_\lambda$. Thus the entire point set and the exclusion are exact.
\end{proof}
This coequalizer is an additional explicitly stated quotient, not a
replacement of ACF19. It records the relation required to make the two
forward maps agree. Its $p$-Frobenius on points is $a\mapsto pa$ on
$\mathbb Z_p$, which is injective and not surjective; in contrast
ACF22 is bijective on the unquotiented full point set. The loss of the
inverse Frobenius is exact: $T^1=1$ does not imply $T^{1/p}=1$.

The support-local semiring need not be Boolean. For the three-element
chain $L=\{0<a<1\}$ and the prime ideal $J=\{0,a\}$, the only
allowed lattice denominator is $1$, so ACF12 gives $L_J=L$ with all
three elements distinct. For $J=\{0\}$, inverting $a$ identifies $a$
with $1$ and leaves two elements. Thus the actual support stalks retain
the different local lattice geometries even though every spherical
stalk in ACF10 over them is $\mathbf F_1$.

The comparison quotient also gives an exact Frobenius calculation.
Under the character parameter its point multiplication is addition in
$\mathbb Z_p$, because $\psi_p((a+b)h)=\psi_p(ah)\psi_p(bh)$.
Writing $m=p^a d>0$ with $(p,d)=1$, its Frobenius therefore has
\begin{equation}\tag{ACF32}
 \begin{gathered}
 \ker(m:\mathbb Z_p\to\mathbb Z_p)=0,\qquad
 \operatorname{coker}(m:\mathbb Z_p\to\mathbb Z_p)
                   =\mathbb Z_p/p^a\mathbb Z_p
                   \simeq\mathbb Z/p^a\mathbb Z,\\
 \{\text{quotient characters}\}\cap\{\text{archimedean local characters}\}
       =\{\chi_{2\pi i k}:k\in\mathbb Z\}.
 \end{gathered}
\end{equation}
The first line follows because $\mathbb Z_p$ is an integral domain and
$d$ is a unit, with the last quotient identified by its first $a$
$p$-adic digits. The second follows from ACF26 and
$\mathbb Z_p\cap\mathbb Z[1/p]=\mathbb Z$. Thus the full algebraic
point group has the prime-to-$p$ kernel $\mu_d$ in ACF22, whereas this
specific comparison quotient has the $p$-primary cokernel in ACF32.
Both statements retain their respective original domains.

\subsection{An exact diagram of the retained data}
For $L=\{0,a,b,1\}$ with $a\wedge b=0$ and $a\vee b=1$, there are two
support primes $P_{\{0,a\}}$ and $P_{\{0,b\}}$, both below $Q_{(0)}$.
The following diagram displays the maps and stalk values; every
horizontal arrow points in the generization direction for stalk maps:
\begin{equation}\tag{ACF29}
\begin{CD}
 \mathbf F_1[T^{H_p^+}] @>{\epsilon_p}>> \mathbf F_1
                              @>{\operatorname{id}}>> \mathbf F_1\\
 @V{T^h\mapsto1}VV @VV{\rm unit}V @VV{\rm unit}V\\
 H G_L(\mathbb Z_{(p)}) @>>> H G_L(\mathbb Q) @>>> H L_J .
\end{CD}
\end{equation}
The three columns lie over $Q_{(p)},Q_{(0)},P_J$, and the last column
is present separately for each of the two support primes. Its bottom
horizontal map is support localization, not an amplitude map.
For this Boolean lattice $L_J\cong\{0,1\}$: localizing the unique
nontrivial element outside $J$ makes it $1$, and its complementary
element becomes $0$. Before this localization the finite-prime point
boundary consists of the four distinct characters
$Z_0,Z_a,Z_b,Z_1$. Amplitude collapses those four to one zero character;
locality selects only the group part; neither operation silently changes
the other. The source locators for the top row are FF Proposition
\texttt{prop12} and its Frobenius structure; the bottom row is proved
in ACF12--ACF16, and the four labelled characters in ACF19--ACF21.


\subsection{The assembled point object and its exact local map}
The base, all its spherical stalks, and all their algebraic characters
can be assembled without any implicit identification. Define the set
of base-labelled algebraic points, with its displayed projection to
$X_L$, by
\begin{equation}\tag{ACF30}
 \begin{split}
 \mathfrak P_L={}&
 \coprod_{J\in\operatorname{Spec}_{\rm lat}L}\{(P_J,*)\}
 \ \sqcup\ \{(Q_{(0)},*)\}\\
 &\sqcup\coprod_{p\ \mathrm{prime}}
       \bigl(\{Q_{(p)}\}\times(\mathcal U_p\sqcup L)\bigr).
 \end{split}
\end{equation}
This is a set with explicitly specified fibre monoids; no additional
topology on this set is asserted. Its structure is fixed by the sheaf
$\mathcal F_L$, the maps ACF11, and the products ACF21. The point map
for a generization from $Q_{(p)}$ to $Q_{(0)}$ or $P_J$ goes in the
opposite direction and selects the constant unit character in the
finite-prime fibre. The multiplicative point at each support prime
retains its base index $J$ even when the corresponding spherical
stalk map is identical to that at another support prime.

Define the place-labelled local object, retaining the trivial character
at every place, by
\begin{equation}\tag{ACF31}
 \begin{gathered}
 \mathfrak P_L^{\rm loc}=
 \coprod_{p\ \mathrm{prime}}
 \left((\{(Q_{(p)},\infty)\}\times\mathbb C)
       \sqcup(\{(Q_{(p)},p)\}\times\mathbb Q_p)\right),\\
 \ell:\mathfrak P_L^{\rm loc}\longrightarrow\mathfrak P_L,\qquad
 (Q_{(p)},\infty,z)\longmapsto(Q_{(p)},\chi_z),\quad
 (Q_{(p)},p,a)\longmapsto(Q_{(p)},\chi_a).
 \end{gathered}
\end{equation}
All fibres of $\ell$ have size at most two. They have size two exactly
at $\chi_{2\pi i b}=\chi_b$ with $b\in\mathbb Z[1/p]$, by ACF26;
their two preimages are the displayed distinct place-labelled points.
The map has no image in the zero-labelled boundary, and there are no
place-labelled points above $P_J$ or $Q_{(0)}$, because their group is
$\{0\}$ and its place set is empty. These statements follow respectively
from ACF19, the unit criterion, ACF23--ACF26, and the source place
definition. Thus ACF30--ACF31 give the full object and its exact local
subfamilies, rather than a collapse to its classical amplitude.

Every positive source Frobenius fixes the base, fixes every boundary
label, and acts on the local parameters by $z\mapsto mz$ and
$a\mapsto ma$. Exponent-zero Frobenius sends every finite-prime
algebraic point to its constant unit character and sends both local
parameters to zero while retaining their place labels. Consequently
$\ell$ is equivariant for every $m\ge0$, directly by ACF22--ACF25.
